
5. 修改后 LaTeX 文本

```tex
% !TEX root = AImath.v5.tex
\chapter*{致谢}
\addcontentsline{toc}{chapter}{致谢}

写作是一项孤独的事业。形单影只，与键盘为伴，上一秒文思泉涌，下一秒却陷入自我怀疑。

但从另一些方面来说，写作又是一种高度社会化的工作。每一位作者都可以回头追问：是谁点燃了我们的创作欲望？是谁在关键时刻给出了那一句提醒？我们又曾从哪些作家、思想家、研究者以及各种不同的思想中汲取了营养。正是这些可见与不可见的同行者，让一本书慢慢长成现在的样子。

无论他们是否知道，下列个人与机构都给予了我们特别的启发与帮助，谨致谢忱。

首先，我们由衷感谢本书的中英文版稿件贡献者和论坛用户们。作为一部开源教材，这本书从一开始就离不开社区的参与：有的读者提交勘误和补充内容，有的读者在论坛上提出尖锐的问题，有的读者则贡献了整段代码或案例。感谢中英文草稿的数百位撰稿人，他们帮助增添或改进了书中的内容，并提供了宝贵的反馈。特别地，我们要感谢这份中文稿的每一位撰稿人，是他们的无私奉献让这本书变得更好。这些贡献者的 GitHub 用户名或姓名是（排名不分先后）：…… 谢谢你们为每一位读者改进这本开源书。

本书的初稿曾在多门课程中被用作教学材料，包括中国科学技术大学、上海财经大学的“深度学习”课程，浙江大学的“物联网与信息处理”课程，以及上海交通大学的“面向视觉识别的卷积神经网络”课程。我们在此感谢这些课程的师生，他们既是本书最早的一批读者，也是最严谨的“测试者”。特别感谢连德富教授、王智教授和罗家佳教授，感谢他们在课堂使用、作业设计与学生反馈方面提出的中肯意见，这些建议直接推动了本书结构与内容的改进。

在写作与实验过程中，我们得到了来自 Amazon Web Services 的大力支持。在此特别感谢 Swami Sivasubramanian、Raju Gulabani、Charlie Bell、Peter DeSantis、Adam Selipsky 和 Andrew Jassy 在本书撰写过程中给予的慷慨支持。如果没有可用的时间、资源，以及来自同事们的讨论和鼓励，这本书的项目就很难真正启动并坚持下来。我们还要感谢 Apache MXNet 团队实现了很多本书所使用的特性，为读者能在实践中复现实验提供了扎实的技术基础。

在内容质量方面，我们受益于许多同事细致耐心的校勘工作。经过他们一轮又一轮的审阅，本书的质量得到了极大的提升。在此我们一一列出章节和校勘人，以表示我们由衷的感谢：引言的校勘人为金颢，预备知识的校勘人为吴俊，深度学习基础的校勘人为张航、王晨光、林海滨，深度学习计算的校勘人为查晟，卷积神经网络的校勘人为张帜、何通，循环神经网络的校勘人为查晟，优化算法的校勘人为郑帅，计算性能的校勘人为郑达、吴俊，计算机视觉的校勘人为解浚源、张帜、何通、张航，自然语言处理的校勘人为王晨光，附录的校勘人为金颢。谢谢你们对每一个符号、每一行公式、每一处表述的严格把关。

我们还要感谢将门创投，特别是王慧、高欣欣、常铭珊和白玉，为本书的两位中国作者在讲授“动手学深度学习”系列课程时所提供的平台和组织工作。感谢所有参与这一系列课程的数千名同学，是你们的提问、作业和反馈，让我们看到了教材在真实课堂中的表现，也帮助我们发现并修补了许多不足之处。感谢 Amazon Web Services 中国团队的同事们，特别是费良宏和王晨，在课程与书稿推进过程中的支持与鼓励。

感谢本书论坛的三位版主：王鑫、夏鲁豫和杨培文。他们牺牲了自己宝贵的休息时间，在论坛上耐心地回复大家的提问、维护讨论秩序，让这本书在纸面之外拥有了一个持续生长的学习空间。感谢人民邮电出版社的杨海玲编辑，在本书出版过程中为我们协调流程、沟通细节、提出专业建议，让这本书最终以更稳定的质量呈现在读者面前。

在具体实现层面，我们还要感谢为不同深度学习框架改编代码的贡献者。感谢 Anirudh Dagar 和唐源，将部分较早版本的 MXNet 实现分别改编为 PyTorch 和 TensorFlow 实现，使更多读者可以在自己熟悉的框架下动手实践。感谢百度团队将较新的 PyTorch 实现改编为 PaddlePaddle 实现，进一步拓展了本书在工业界与国产框架生态中的适用性。感谢张帅将更新的 LaTeX 样式集成进 PDF 文件的编译流程，让本书以更清晰、美观的排版呈现出来。

最后，我们要特别感谢我们的家人。谢谢你们一直默默陪伴着我们，在我们埋头写作、反复修改、偶尔沮丧的日子里，承担了许多看不见的时间与情绪成本。没有你们的理解与支持，这本书无法按现在的样子与大家见面。

\newpage
```
